\section{Pavyzdinės programos}

\subsection{Amžinas ciklas}

\begin{verbatim}
    ; Duomenys apie failą šešioliktaine:
    494E4C4F4F50 4558 ; Failo pavadinimas ir failo tipas - INLOOP EX
    00000008 ; Failo dydis baitais
    00000000 ; Sekančio segmento adresas
    00000008 ; Dabartinio segmento dydis
    
    00000000 ; Pirmos instrukcijos žodžio adresas
    
    ; Instrukcijos ASCII:
    JP00 ; Grįžtam į pirmą instrukciją
\end{verbatim}

Visas programos failas šešioliktaine:

49 4E 4C 4F 4F 50 45 58 00 00 00 08 00 00 00 00 00 00 00 08 00 00 00 00 4A 50 4F 4F

\subsection{Įvestis/išvestis}

\begin{verbatim}
    ; Duomenys apie failą šešioliktaine:
    494E504F5554 4558 ; Failo pavadinimas ir failo tipas - INPOUT EX
    0000003C ; Failo dydis baitais
    00000000 ; Sekančio segmento adresas
    0000003C ; Dabartinio segmento dydis
    
    00000002 ; Pirmos instrukcijos žodžio adresas

    ; Programos duomenys šešioliktaine:
    00000001 00000000 01000000 ; Nustatome 3 žodžius su duomenimis 
                      ; (simbolių sk. ir vienas žodis saugoti gautai reikšmei + skaičius sudėčiai)
                      
    ; Instrukcijos ASCII:
    LD00 ; Į R1 įkeliame simbolių skaičių
    IN01 ; Skaitome simbolį ir įrašome į atmintį adresu 01
    LD01 ; Į R1 įkeliame nuskaitytą simbolį
    SWAP ; Apkeičiame R1 ir R2
    LD02 ; Į R1 įkeliame skaičių sudėčiai
    ADDR ; Sudedame registrus (simbolį padidiname vienetu)
    SWAP
    ST01 ; Išsaugome suskaičiuotą reikšmę
    LD00 ; Į R1 įkeliame simbolių skaičių
    OU01 ; Spausdiname ekrane
    HALT ; Baigiame programą
\end{verbatim}

Visas programos failas šešioliktaine:

49 4E 50 4F 55 54 45 58 00 00 00 3C 00 00 00 00 00 00 00 3C 00 00 00 02 00 00 00 01 00 00 00 00 01 00 00 00 4C 44 30 30 49 4E 30 31 4C 44 30 31 53 57 41 50 4C 44 30 32 41 44 44 52 53 57 41 50 53 54 30 31 4C 44 30 30 4F 55 30 31 48 41 4C 54

\subsection{Failų skaitymas/įrašymas (nebaigta)}
\begin{verbatim}
    ; Duomenų failas:
    545854444154 5458 ; Failo pavadinimas ir failo tipas - TXTDAT TX
    00000010 ; Failo dydis baitais
    00000000 ; Sekančio segmento adresas
    00000010 ; Dabartinio segmento dydis

    ; Failo duomenys šešioliktaine:
    00000000 00000000 00000000 00000000

    ; Duomenys apie failą šešioliktaine:
    465445535453 4558 ; Failo pavadinimas ir failo tipas - FTESTS EX
    00000064 ; Failo dydis baitais
    00000000 ; Sekančio segmento adresas
    00000064 ; Dabartinio segmento dydis
    
    00000008 ; Pirmos instrukcijos žodžio adresas

    ; Programos duomenys šešioliktaine:
    54585444 41545458 00000000; Atidaromo failo pavadinimas ir žodis failo numeriui
    68656C6C 6F000000 ; Žinutė, kurią rašysime į failą ("hello")
    00000000 00000000 ; Kur rašysime skaitomus duomenis
    00000005 ; Rašomų/skaitomų simbolių skaičius
    ; Instrukcijos ASCII:
    
    FO00 ; Atidaromas failas
    ST02 ; Išsaugome failo numerį
    LD07 ; Į R1 rašome rašomų simbolių skaičių
    SWAP ; Apkeičiame R1 ir R2
    LD02 ; Į R1 įsikeliame failo numerį
    PD03 ; Rašome į failą
    FCLR ; Uždarome failą
    FO00 ; Atidaromas failas
    ST02 ; Išsaugome failo numerį
    LD07 ; Į R1 rašome skaitomų simbolių skaičių
    SWAP ; Apkeičiame R1 ir R2
    LD02 ; Į R1 įsikeliame failo numerį
    GD05 ; iš failo rašome į atmintį
    FCLR ; Uždarome failą
    LD07 ; Į R1 rašome išvedamų simbolių skaičių
    OU05 ; iš atminties išvedame į ekraną
    HALT ; baigiama programa
\end{verbatim}

Visas duomenų ir programos failas šešioliktaine:

54 58 54 44 41 54 54 58 00 00 00 10 00 00 00 00 00 00 00 10 00 00 00 00 00 00 00 00 00 00 00 00 00 00 00 00 46 54 45 53 54 53 45 58 00 00 00 64 00 00 00 00 00 00 00 64 00 00 00 08 54 58 54 44 41 54 54 58 00 00 00 00 68 65 6C 6C 6F 00 00 00 00 00 00 00 00 00 00 00 00 00 00 05 46 4F 30 30 53 54 30 32 4C 44 30 37 53 57 41 50 4C 44 30 32 50 44 30 33 46 43 4C 52 46 4F 30 30 53 54 30 32 4C 44 30 37 53 57 41 50 4C 44 30 32 47 44 30 35 46 43 4C 52 4C 44 30 37 4F 55 30 35 48 41 4C 54